\input{../tex-inputs/latex-leseansicht-vorspann}

\section[Gustav Schwarzkopf an Arthur Schnitzler, 3. 7. 1914]{L04249 Gustav Schwarzkopf an Arthur Schnitzler, 3. 7. 1914}
\nopagebreak\mylabel{L04249v}
\rehead{ }\normalsize\beginnumbering\briefempfaengerindex{Schnitzler, Arthur@\textsc{Schnitzler, Arthur}!zzzSchwarzkopf, Gustav@\emph{von Gustav Schwarzkopf}!1914-07-031@{3. 7. 1914}|(be}
\toendnotes[C]{\smallbreak\pagebreak[2]}
\correspDesc{Versand  durch Gustav Schwarzkopf am 3. 7. 1914 in Weißenbach am Attersee
\newline{}Erhalt  durch Arthur Schnitzler im Zeitraum [4. 7. 1914 – 8. 7. 1914?] in Wien}\toendnotes[C]{\smallbreak}
\Standort{CUL, Schnitzler, B 96a.}
\physDesc{Bildpostkarte, 255 Zeichen
\newline{}Handschrift: schwarze Tinte, deutsche Kurrent
\newline{}Versand: Stempel: »\nobreak{}\oindex{Weißenbach am Attersee@\textbf{Weißenbach am Attersee}, \emph{Verwaltungsgebiet}|pwk}Weissenbach am Attersee, 3. VII. 14\nobreak{}«.  
\newline{}Schnitzler: mit Bleistift Vermerk: »\textsc{Schwarzkopf}« }\toendnotes[C]{\smallbreak}\pstart{}{\pb}Herrn \textsc{D\textsuperscript{r} Arthur Schnitzler}\pend{}\pstart{}\textsc{Wien
                  XVIII\oindex{XVIII., Währing@\textbf{XVIII., Währing}, \emph{Verwaltungsgebiet}|pw}}\pend{}\pstart{}Sternwarteſtraſſe 71\oindex{Wien@\textbf{Wien}!XVIII., Währing@\textbf{XVIII., Währing}!Sternwartestraße 71@\textbf{Sternwartestraße 71}, \emph{Wohngebäude}|pw}\pend{}{\bigskip}
\pstart
           \noindent{}\centering{}{\pb}\textcolor{gray}{\textbf{Salzkammergut. Weissenbach am Altersee}}\oindex{Weißenbach am Attersee@\textbf{Weißenbach am Attersee}, \emph{Verwaltungsgebiet}|pw}\pend
           \vspace{1em}
\pstart
           \raggedleft{}{\pb}3. VII. 14.\pend
           \vspace{0.5em}
\pstart
           Zu meiner Überraſchung bin ich{ }ſeit geſtern Nachmittag wirklich hier; war heute{ }ſogar ſchon in \label{K_L04249-1v}\edtext{\textsc{Unterach\oindex{Unterach am Attersee@\textbf{Unterach am Attersee}|pw}} und habe dort Herrn\pwindex{Salten, Felix 6.\,9.\,1869 Budapest – 8.\,10.\,1945 Zürich@\textsc{Salten, Felix} (6.\,9.\,1869 Budapest – 8.\,10.\,1945 Zürich), \emph{Schriftsteller, Journalist, Chefredakteur}|pw} und Frau Salten\pwindex{Salten, Ottilie 7.\,3.\,1868 Prag – 22.\,6.\,1942 Zürich@\textsc{Salten, Ottilie} (7.\,3.\,1868 Prag – 22.\,6.\,1942 Zürich), \emph{Schauspielerin}|pw}}{\lemma{\textnormal{\emph{Unterach … Salten}}}\Cendnote{\textnormal{Diese
                  verbrachten den Sommer im Berghof\oindex{Berghof@\textbf{Berghof}, \emph{Wohngebäude}|pwk} in Unterach am Attersee\oindex{Unterach am Attersee@\textbf{Unterach am Attersee}|pwk}, vgl. XXXX Auszeichnungsfehler: Dokument L03565 nicht gefunden. }}}\label{K_L04249-1} getroffen.\pend
           
\pstart
           Ich grüße ihre Frau\pwindex{Schnitzler, Olga 17.\,1.\,1882 Wien – 13.\,1.\,1970 Lugano@\textsc{Schnitzler, Olga} (17.\,1.\,1882 Wien – 13.\,1.\,1970 Lugano), \emph{Schauspielerin, Sängerin}|pwv}
               und{\\[\baselineskip]}Sie herzlichſt.{\\[\baselineskip]}Ihr{\\[\baselineskip]}\spacefill\mbox{Gustav.}\pend
           \leftskip=0em{}\selectlanguage{ngerman}\endnumbering\briefempfaengerindex{Schnitzler, Arthur@\textsc{Schnitzler, Arthur}!zzzSchwarzkopf, Gustav@\emph{von Gustav Schwarzkopf}!1914-07-031@{3. 7. 1914}|)be}\mylabel{L04249h}
\begin{anhang}
\end{anhang}\newcommand{\dateiname}{L04249}\newcommand{\titel}{Gustav Schwarzkopf an Arthur Schnitzler, 3. 7. 1914}\newcommand{\editorInnen}{Herausgegeben von Jahnke, SelmaMüller, Martin Anton}\input{../tex-inputs/latex-leseansicht-abspann}
